% !TeX root = ../papers/nc-p3.tex
\section{任意の成分における有界PI次数のNoether窓の有限判定}
\label{sec:noetherian-pi-window-cutoff}
\begin{theorem}\label{thm:noetherian-pi-window-cutoff}
正整数 $s$ を固定し、$h=4\operatorname{lcm}(1,\ldots,s)\operatorname{lcm}(1,\ldots,s^2)$
とする。任意の標準有限Frobenius候補について次は同値である。
\begin{enumerate}
\item 標準Koszul左右AS正則で、標準折り畳み $B$ が左右Noetherかつ
$4s$ 次行列環の多項式恒等式を満たす。
\item $B$ に経路長 $h$ の中心四元が存在して定理\ref{thm:successively-central-cutoff}の
有限検査を通り、得られた四元上の有限自由代数の一般商のPI次数が $4s$ 以下である。
\end{enumerate}
成分、点対応の有限性、一般窓からの変形を仮定しない。検査に入る窓は全て合成整域で、
標準折り畳みは正規な有限生成中心上有限となる。とくに十五次元閉包の有限四辺では
定理\ref{thm:finite-domain-specialization-parameters}の564次検査が正則かつ左右Noetherであるためにも
必要十分である。
\end{theorem}
\begin{proof}[成長の小さい部分加群の排除]
まず左右NoetherでPIな正則窓を考える。有限生成次数付き右加群には有限生成頂点射影による有限分解があり、
各Hilbert級数は $p(t)/(1-t)^4$、$p\in\mathbb Z[t,t^{-1}]$ である。
GK次元が $d$ 以下のそのような加群 $M$ は $j<4-d$ で $\underline{\operatorname{Ext}}^j_B(M,B)=0$ を満たす。
$d=0$ は有限長のAS消滅である。一般に $M$ には、各因子 $N$ の全ての非零な斉次部分加群が同じ
両側消去イデアル $P$ を持つ有限フィルトレーションがある。各残余商で非零部分加群の消去イデアルを
極大に選んでその部分加群を加え、昇鎖条件を使う。$P$ は素である。実際 $IJ\subset P$、$I\not\subset P$
なら $NI\ne0$ なので $J\subset P$ である。まず斉次イデアルについて従い、最高斉次項で通常の素性も従う。
$B/P$ が無限次元ならその正次数イデアルは非零の斉次中心元 $z$ を含む。
素PI環の中心局所化が単純Artin環であり、局所化して単位元を表示した後に中心分母を払えばよい。
中心元の各斉次部分は再び中心元である。$z$ は $N$ に単射に作用する。核が非零なら核の消去は
$P$ より真に大きくなる。余核のGK次元はHilbert級数から少なくとも一つ落ちる。
帰納仮定と長完全列により $j<4-d$ で $\underline{\operatorname{Ext}}^j_B(N,B)$ は正次数移動と同型で、
有限生成射影分解から内部次数が下に有界なので零である。$B/P$ が有限次元の因子は既に有限長である。
フィルトレーションで主張が従う。

よって $e_iB$ の非零部分加群はGK次元4を持ち、その分子の $p(1)$ は正整数である。
$e_iB$ 自体ではこの整数が1なので、非零な斉次経路元による頂点射影間の写像では像が1を使い切り、
核は零である。最高次を取ると全ての非零合成が非零になる。$B$ は素で、商環は $Q=M_4(D)$、
$D$ は頂点隅の商可除環である。実際各原始頂点冪等元は中心局所化後もその隅が整域なので原始である。
この $D$ の中心上の次数を $a$ とすると $a\le s$ である。
\end{proof}

\begin{lemma}\label{lem:cy-finite-birational-transfer}
$B$ を上の標準正則窓の素な折り畳みとする。$Q$ を中心局所化して得る単純Artin環とし、
$B\subset T\subset Q$ を次数と両立する環拡大とする（非斉次分母も入れた $Q$ における斉次分数を使う）。
$T$ が左右それぞれ $B$ 上有限生成次数付き射影加群による有限分解を持てば $Z(T)\subset Z(B)$ である。
\end{lemma}
\begin{proof}
両側Koszul分解の双対は最終次数4に $B(4)$ だけを持つ。最終の符号が内部共役で消える計算は
定理\ref{thm:central-cut-center-dichotomy}の相対双対の証明の通りで、そこでは中心列を使わない。
左と右で完全な両側加群 $M$ の二つの導来双対 $M^r=\operatorname{RHom}_{\mathrm{mod}\,B}(M,B)$ と
$M^l=\operatorname{RHom}_{B\,\mathrm{mod}}(M,B)$ の間に両側加群として自然な同型がある。
実際 $P_\bullet$ を $B$ の上の有限両側射影分解とし、$P_\bullet\otimes_B M$ の両側双対を取ると
$M^r\otimes_B^{\mathbf L}\operatorname{RHom}_{B^e}(B,B^e)$ である。
各自由項 $B\otimes W\otimes B$ では右 $B$ 完全性によるテンソル・Hom同型であり、
微分の両側乗法とも両立する。$M\otimes_B P_\bullet$ から同じ両側双対を計算すれば
$\operatorname{RHom}_{B^e}(B,B^e)\otimes_B^{\mathbf L} M^l$ となる。
両結果から共通の内部および複体移動を外す。同型 $M^r\simeq M^l$ は $M=B$ で恒等となるよう
可逆な中心スカラーで正規化できる。

$z\in Z(T)$ の作用を $m_z$ とする。右側の余評価、上の同型、左側の評価の合成
\[
 B\longrightarrow T\otimes_B^{\mathbf L} T^r
 \xrightarrow{m_z\otimes 1} T\otimes_B^{\mathbf L} T^l\longrightarrow B
\]
は $B$ の導来両側自己射なのでその次数零コホモロジーで中心元を与える（内部移動は $z$ の次数）。
余評価は恒等写像、評価は左線型汎関数の値である。中心の分母を逆にすると $B\hookrightarrow T$ は恒等
$Q\to Q$ になり、有限射影双対もこの局所化と可換である。双対間の同型は包含に自然であり、従って
局所化された合成は $Q$ 上の $z$ 倍そのものである。$B\hookrightarrow Q$ の単射から従う。
\end{proof}

\begin{proof}[被約traceの有限拡大と判定次数]
$Q$ の被約traceを $\tau$ とし、$T$ を $B$ にその全元の被約traceを添加した部分環とする。
これは次数と両立する。$T$ は有限生成可換部分環上有限である。具体的には有限生成元を分裂体上の
$4a$ 次行列として表し、同じ個数とサイズの一般行列から特殊化する。
一般行列のtrace環はその有限生成のスカラー不変式環上有限である（標数零の行列不変式論）。
これは同時共役不変式が語のtraceで生成され、行列値同変式が不変式上有限加群となることからも従う。
この有限性は像にも移り、頂点射影も生成元に含められる。

さらに $T$ は $B$ の左右どちらでも有限生成である。そのための分母を説明する。
$n=(4a)^2$ とし、$n$ 個の $x$ を交代化した積 $x_1 y_1 x_2\cdots y_{n-1}x_n$ を $F(x,y)$ と書く。
$Q$ でこれは恒等的に零ではない。分裂後に全行列単位を $x$ に並べ、隣り合う単位をつなぐ
行列単位を $y$ に入れると、最初の行と最後の列を固定した成分には指定の順だけが寄与する。
多重線型性と中心局所化から $B$ 内にも非零な値 $f$ がある。
交代な最上外冪で左 $b$ 倍の導分作用を取れば
\[
 (4a)\tau(b)F(x,y)=\sum_{j=1}^n F(x_1,\ldots,bx_j,\ldots,x_n;y).
\]
従って任意個のtraceの積を $f$ に掛けても $B$ に入り、$fT\subset B$。
この性質を持つ $B$ の元全体は非零両側イデアルで、中心局所化で単位を生成するため非零な中心分母 $c$ を含む。
$cT\subset B$ は左右イデアルで、左右Noether性からどちらでも有限生成となる。よって $T$ は左右完全であり、
補題\ref{lem:cy-finite-birational-transfer}により添加したtraceは全て $Z(B)$ に属する。
これで $B=T$ はその有限生成中心上有限である。中心自体も正規である。
実際中心分数体内の整元の斉次整閉包を添加した環にも同じ補題が適用できる。
有限型次数付き整域の正規化は有限かつ次数付きで、$B$ も既にその中心上有限である。

各非頂点中心指標で単純加群 $S$ を取る。各頂点次元は同じ $d\le a$ である。
非零な正次数中心値で展開した周期的加群は有限生成で、そのHilbert級数は
$(\sum_{i=0}^3\dim(Se_i)t^i)/(1-t^4)=p(t)/(1-t)^4$ だからである。上界は隅のPI次数による。
固定した頂点で長さ4のループの像は単位と共に $M_d(k)$ を生成し、少なくとも一つのループ $X$ は零でない。
実際隅の単純作用にBurnsideの定理を使え、全ループの零作用は正次数中心値の存在に反する。
単位から始め、長さ4の生成元を掛けた累積張空間は安定すれば全環である。ゆえに高々 $d^2-1$ 個の積で
$M_d$ が張られ、trace対合から $\operatorname{tr}(XY)\ne0$ となるそのような語 $Y$ がある。
$b=XY$ はこの単純上非冪零で、経路長 $4j$ は $1\le j\le s^2$ を満たす。
一般隅の被約特性係数は $B$ の全中心に取れる。実際零番以外の対角を零にした $b$ に対する
全行列の被約traceは隅の被約traceそのもので、冪にも同じことを適用しNewton公式を使う。
次数 $4ji$ ($1\le i\le a$) のこの係数が全て指定指標で零なら $b^a$ はその隅で零作用になり矛盾する。
各 $4ji$ が $h$ を割るので $Z(B)_h$ は射影中心上で基点がない。三次元の射影中心を四一般元で切れば
剰余は有限次元になり、必要な検査を得る。逆に検査から正則性と中心四元上の有限自由性が従い、
一般商の次数は有限な乗法表の中心を求めてその上の次元の平方根を取ればよい。
素性は頂点射影の重複度1による既出の証明による。最後の閉包上では $B$ は常に次数16の行列恒等式を満たし、
一隅整域なら既出の短い検査が使える。
\end{proof}
